\section{Task Modality and Benchmark (Carrier)}
\label{sec:benchmark}
% 【中文】\section{任务模态与基准（载体）}

\bench{} is a benchmark of \textbf{\Ntasks{} tasks} across \textbf{13 EDA domains}---from schematic/netlist to SPICE simulation, physical verification, and digital synthesis---each derived from \emph{real coding-agent traces} of IC design engineers in a production environment. We use \bench{} as a \emph{carrier} for falsifiable cognition claims, not as the contribution itself.
% 【中文】
% \bench{} 是一个含 \textbf{\Ntasks{} 个任务}、跨越 \textbf{13 个 EDA 领域}的基准——从原理图/网表到
% SPICE 仿真、物理验证与数字综合——每个任务源自生产环境中 IC 设计工程师的\emph{真实编程智能体轨迹}。
% 我们将 \bench{} 用作可证伪认知主张的\emph{载体}，而非贡献本身。

\textbf{Design-flow contract.} Each task requires the agent to produce \texttt{submit.zip} by orchestrating remote EDA tools through an MCP bridge, scored against $\geq 8$ executable assertions with a \texttt{PreToolUse} hook that prevents answer leakage (reference/ and check.py never enter the workdir). The contract deliberately does not specify language or host---the agent must discover the right tool chain. Detailed benchmark construction is in Appendix~\ref{sec:benchmark-details}.
% 【中文】
% \textbf{设计流契约。} 每个任务要求智能体通过 MCP 网桥编排远程 EDA 工具产出 \texttt{submit.zip}，
% 按 $\geq 8$ 条可执行断言评分，并用 \texttt{PreToolUse} 钩子防止答案泄漏（reference/ 与 check.py
% 永不进入工作目录）。该契约故意不规定语言或主机——智能体必须自行发现正确的工具链。详细基准构建见
% 附录~\ref{sec:benchmark-details}。

\textbf{Evaluation subset.} Experiments use a Full-set subset (top-120 by grounding depth / chain length, internally evaluated, not released). A frozen 696-run snapshot (2026-11-09; DeepSeek-V4-Pro, Cline v3.0.51, \texttt{thinking=low}, 2400s timeout, 8 workers) provides the M0 baseline anchor: \MzeroPass{} Pass@1 (62 PASS / 558 FAIL / 70 ERROR / 6 ABSENT).
% 【中文】
% \textbf{评测子集。} 实验使用 Full 集子集（按接地深度/链长取 top-120，内部评测，不发布）。一份冻结
% 的 696-run 快照（2026-11-09；DeepSeek-V4-Pro、Cline v3.0.51、\texttt{thinking=low}、2400s 超时、
% 8 workers）提供 M0 基线锚点：\MzeroPass{} Pass@1（62 PASS / 558 FAIL / 70 ERROR / 6 ABSENT）。

\textbf{Open-source boundary (v1.4 narrowing).} We release a \emph{protocol white-paper} (harness protocol + criteria + desensitized task specifications) and an \emph{NDA-controlled review package} for reviewers, but do \textbf{not} release executable Lite-set assets or a mock-EDA mode. The falsifiability of our claims does not depend on releasing proprietary EDA tool licenses: the protocol and criteria are sufficient for an independent party to re-derive an equivalent benchmark on open-source EDA tools (Yosys/OpenROAD/Magic) and test the same hypotheses (\S\ref{sec:reproducibility}).
% 【中文】
% \textbf{开源边界（v1.4 收缩）。} 我们发布\emph{协议白皮书}（执行框架协议 + 判据 + 脱敏任务规格）与
% \emph{NDA 受控评审包}供评审者使用，但\textbf{不}发布可执行的 Lite 集资产或 mock-EDA 模式。我们主张的
% 可证伪性不依赖于发布专有 EDA 工具许可：协议与判据足以让独立方在开源 EDA 工具（Yosys/OpenROAD/Magic）
% 上重新推导出等价基准并检验同样假设（\S\ref{sec:reproducibility}）。
